\section{From FTRFS to beamfs}\label{sec:ftrfs}

FTRFS~\cite{fuchs2015ftrfs} proposed a filesystem for the on-board
memory of spacecraft: checksums, forward error correction and memory
protection over a small radiation-exposed store. Its Linux realisation,
posted as an RFC in April 2026~\cite{ftrfsrfc} and described in the
FTRFS v1 report~\cite{desbrieres2026ftrfs}, inherited the concept and
its limits. Each limit below is conceptual, not a missing feature: it
follows from the threat model or from the place error correction holds
in the design, which is why \bfs is a new filesystem with its own
on-disk format rather than a further FTRFS revision.
\Cref{tab:ftrfs} sets them side by side.

\paragraph*{Threat model}
FTRFS assumes radiation-induced upsets arriving as independent events at
a low constant rate, the Poisson regime of space radiation. It says
nothing about correlated events, bursts, or a deliberate attacker. The
soundness theorem of the first \bfs report was stated under that model;
fault injection falsified it, and v2 withdrew it~\cite{desbrieres2026beamfsv2}.
\bfs states two families of perturbation, claims one, and says which
(\cref{sec:scope}).

\paragraph*{Where correction sits}
In FTRFS, Reed-Solomon protects metadata and a CRC guards each block;
the decoder runs once a CRC has failed, and the RFC shipped the encoder
only~\cite{ftrfsrfc}. File data, the part an application reads back,
had no correction at all, and a CRC that is itself hit can mask or
invent a failure. In \bfs every data block is encoded, and the decoder
runs on every read whatever any checksum says (\cref{sec:readpath}).

\paragraph*{Scale}
The FTRFS allocation bitmap is one RS-protected block: 16 codewords of
239 bytes, 30\,592 bits, so at most 30\,592 blocks (about 119\,MiB at
4\,KiB). The \bfs superblock records up to 65\,535 bitmap blocks.

\paragraph*{Kernel integration}
The RFC had no \texttt{address\_space\_operations}, and its review
recalled the strong aversion to merging new filesystems built on buffer
heads~\cite{ftrfsrfc}. \bfs reads and writes file data through iomap
(\texttt{iomap\_read\_folio}, \texttt{iomap\_readahead},
\texttt{iomap\_file\_buffered\_write}), but metadata (indirect blocks,
bitmap, scrubber) still goes through buffer heads. That objection is
narrowed, not lifted.

\paragraph*{Evidence}
The RFC was tested by mounting a volume and listing an empty directory.
\bfs is tested with the xfstests \texttt{auto} group on two
architectures and under fault injection (\cref{sec:results}).

\begin{table}[t]
\centering
\caption{FTRFS and beamfs on the points that separate them.}
\label{tab:ftrfs}
\footnotesize
\begin{tabularx}{\columnwidth}{@{}>{\raggedright\arraybackslash}p{1.45cm}
  >{\raggedright\arraybackslash}X >{\raggedright\arraybackslash}X@{}}
\toprule
 & FTRFS (2015 concept, 2026 RFC) & \bfs 0.1.26 \\
\midrule
threat model & radiation only, independent upsets &
  two families; B explicitly not claimed \\
file data & CRC, no correction & 16 RS codewords per block \\
decoder & after a CRC failure; encoder only in the RFC &
  on every read \\
bursts & not addressed & data interleaved: a burst must exceed
  128 bytes; parity grouped: 9 bytes suffice (\cref{sec:capsule}) \\
bitmap & one block, 30\,592 blocks max & up to 65\,535 bitmap blocks \\
VFS & no aops; buffer heads & iomap for data; buffer heads for metadata \\
evidence & mount and empty listing & xfstests, two architectures;
  fault injection \\
\bottomrule
\end{tabularx}
\end{table}
